% Chapter 1: Potential Outcomes and the FWL Theorem
% ============================================================================
% This chapter establishes the theoretical foundation for Double Machine Learning:
%   1. Potential outcomes framework for defining causal effects
%   2. The fundamental problem of causal inference
%   3. The Frisch-Waugh-Lovell theorem as motivation for DML
% ============================================================================

\chapter{Potential Outcomes and the Frisch-Waugh-Lovell Theorem}
\label{ch:potential_outcomes_fwl}

\section{Introduction}\label{sec:ch01_introduction}

Causal inference is the science of determining whether and to what
extent one variable influences another. Unlike prediction tasks, where
we seek to forecast outcomes given features, causal inference aims to
answer \textbf{what-if} questions: What would happen if we changed a
treatment assignment? What is the effect of a policy intervention?

This chapter lays the theoretical foundation for Double Machine Learning
by developing three core concepts:

\begin{enumerate}
\def\labelenumi{\arabic{enumi}.}
\item
  \textbf{The Potential Outcomes Framework}
  \cite{rubin1974estimating,holland1986statistics}: A rigorous
  mathematical language for defining causal effects
\item
  \textbf{The Fundamental Problem of Causal Inference}: Why causal
  inference is inherently a missing data problem
\item
  \textbf{The Frisch-Waugh-Lovell Theorem}
  \cite{frisch1933partial,lovell1963seasonal}: A classical
  regression result that motivates the DML approach
\end{enumerate}

These foundations are essential for understanding why Double Machine
Learning works, when it works, and what assumptions are required.

\subsection{Motivating Example: Insurance
Pricing}\label{_motivating_example_insurance_pricing}

Consider a practical problem from actuarial science: An insurance
company wants to understand the causal effect of competitor pricing on
their own sales. The company observes:

\begin{itemize}
\item
  \(Y_i\): Their sales in week \(i\)
\item
  \(T_i\): Average competitor price in week \(i\)
\item
  \(X_i\): Macroeconomic conditions (VIX, consumer sentiment, interest
  rates)
\end{itemize}

A naive regression of \(Y_i\) on \(T_i\) confounds the causal effect
with:

\begin{itemize}
\item
  \textbf{Selection bias}: Competitors may raise prices during
  high-demand periods
\item
  \textbf{Confounding}: Economic conditions affect both competitor
  prices and sales
\item
  \textbf{Reverse causality}: High sales might induce competitors to
  adjust pricing
\end{itemize}

The potential outcomes framework provides a precise language for
defining the causal effect we seek, and the Frisch-Waugh-Lovell theorem
suggests how to isolate it.

\section{The Potential Outcomes
Framework}\label{_the_potential_outcomes_framework}

\subsection{Definition}\label{_definition}

The potential outcomes framework, developed by Neyman (1923) for
experiments and extended by Rubin \cite{rubin1974estimating} to
observational studies, defines causality through counterfactuals.

\begin{definition}[Potential Outcomes]
\label{def:potential_outcomes}
For each unit $i$ and treatment value $t \in \{0, 1\}$, define:
\[
Y_i(t) = \text{Outcome unit } i \text{ would have if assigned treatment } t
\]

Key properties:
\begin{itemize}
\item $Y_i(0)$: \textbf{Potential outcome under control} (what would happen without treatment)
\item $Y_i(1)$: \textbf{Potential outcome under treatment} (what would happen with treatment)
\item Only one is observed: $Y_i = T_i \cdot Y_i(1) + (1 - T_i) \cdot Y_i(0)$
\end{itemize}
\end{definition}

\begin{example}[Insurance Pricing]
\label{ex:insurance_pricing}
For week $i$:
\begin{itemize}
\item $Y_i(1)$: Sales if competitor average price is high ($T_i = 1$)
\item $Y_i(0)$: Sales if competitor average price is low ($T_i = 0$)
\end{itemize}

The potential outcome $Y_i(0)$ is \textbf{counterfactual} when $T_i = 1$: we observe high competitor prices but want to know what sales would have been with low prices.
\end{example}

\subsection{Individual Treatment
Effect}\label{_individual_treatment_effect}

\begin{definition}[Individual Treatment Effect]
\label{def:individual_treatment_effect}
The causal effect of treatment on unit $i$ is:
\[
\tau_i = Y_i(1) - Y_i(0)
\]

This is the difference in outcomes under treatment versus control \textbf{for the same unit}.

\textbf{Key insight}: $\tau_i$ is a \textbf{deterministic quantity}---it is fixed for unit $i$. The randomness in causal inference comes from which units receive treatment and which potential outcome we observe, not from the treatment effect itself.
\end{definition}

\section{The Fundamental Problem of Causal
Inference}\label{_the_fundamental_problem_of_causal_inference}

\textbf{Fundamental Problem} \cite{holland1986statistics}:
\textbf{We can never observe both potential outcomes for the same unit
at the same time.}

For unit \(i\): * If \(T_i = 1\), we observe \(Y_i(1)\) but not
\(Y_i(0)\) * If \(T_i = 0\), we observe \(Y_i(0)\) but not \(Y_i(1)\)

Therefore, the individual treatment effect \(\tau_i = Y_i(1) - Y_i(0)\)
is \textbf{never directly observable}.

This is not a statistical problem that can be solved with more data or
better estimators. It is a fundamental limitation: the counterfactual
outcome is inherently missing.

\subsection{Observed vs. Unobserved}\label{_observed_vs_unobserved}

Let \(Y_i^{\text{obs}}\) denote the observed outcome:

\[Y_i^{\text{obs}} = \begin{cases}
Y_i(1) & \text{if } T_i = 1 \\
Y_i(0) & \text{if } T_i = 0
\end{cases} = T_i Y_i(1) + (1 - T_i) Y_i(0)\]

The counterfactual outcome \(Y_i^{\text{mis}}\) is unobserved:

\[Y_i^{\text{mis}} = \begin{cases}
Y_i(0) & \text{if } T_i = 1 \\
Y_i(1) & \text{if } T_i = 0
\end{cases}\]

\textbf{Implication}: Causal inference is fundamentally a
\textbf{missing data problem}. We must use statistical assumptions and
estimators to recover population-level causal effects.

\section{Average Treatment Effect}\label{_average_treatment_effect}

Since individual treatment effects \(\tau_i\) are unobservable, we focus
on \textbf{population-level} average effects.

\subsection{Definition}\label{_definition_2}

\begin{definition}[Average Treatment Effect]
\label{def:ate}
The Average Treatment Effect (ATE) is the expected difference in potential outcomes:
\[
\text{ATE} = \E[Y_i(1) - Y_i(0)] = \E[Y_i(1)] - \E[Y_i(0)]
\]

This is the average causal effect across the entire population. Unlike $\tau_i$, the ATE \textbf{can} be estimated under appropriate assumptions.
\end{definition}

\subsection{Why the ATE is
Identifiable}\label{_why_the_ate_is_identifiable}

While individual effects \(\tau_i\) are never observed, the ATE can be
estimated by comparing treated and control groups:

\begin{align}
\text{ATE} &= \mathbb{E}[Y_i(1)] - \mathbb{E}[Y_i(0)] \\
&= \mathbb{E}[Y_i | T_i = 1] - \mathbb{E}[Y_i | T_i = 0] \quad \text{(under randomization)}
\end{align}

The second line holds \textbf{if and only if} treatment assignment is
randomized (or "as-if" randomized after conditioning on confounders).

\textbf{Intuition}: In a randomized experiment: *
\(\mathbb{E}[Y_i(1) | T_i = 1\) =
\textbackslash mathbb\{E\}{[}Y\_i(1){]}{]}: Treated units are
representative of the population * \(\mathbb{E}[Y_i(0) | T_i = 0\) =
\textbackslash mathbb\{E\}{[}Y\_i(0){]}{]}: Control units are
representative of the population

Therefore, comparing treated vs. control groups recovers the ATE.

\section{Identification: From Causal Estimands to Statistical
Estimands}\label{_identification_from_causal_estimands_to_statistical_estimands}

Identification asks: \textbf{Under what assumptions can we express a
causal quantity (like ATE) as a function of the observed data
distribution?}

\subsection{Conditional Independence
Assumption}\label{_conditional_independence_assumption}

\textbf{Assumption 1.1 (Unconfoundedness / Conditional Independence)}

Treatment assignment is independent of potential outcomes, conditional
on observed covariates \(X\):

\[\{Y_i(0), Y_i(1)\} \perp\!\!\!\perp T_i \mid X_i\]

\textbf{Interpretation}: After conditioning on \(X\), treatment
assignment is "as-if randomized"---it does not depend on the potential
outcomes.

\textbf{Example}: In our insurance pricing example: * \(X_i\) includes
macroeconomic conditions (VIX, interest rates, consumer sentiment) *
Assumption: Conditional on these, competitor pricing is unrelated to
potential sales outcomes

This is a \textbf{strong assumption} and must be justified by domain
knowledge.

\subsection{Overlap (Positivity)}\label{_overlap_positivity}

\textbf{Assumption 1.2 (Overlap / Positivity)}

For all \(x\) in the support of \(X\):

\[0 < \mathbb{P}(T_i = 1 | X_i = x) < 1\]

\textbf{Interpretation}: Every unit has a positive probability of
receiving treatment and control, regardless of covariates. Without
overlap: * Some covariate values are only observed in treated group $\to$
Cannot estimate \(\mathbb{E}[Y(0) | X = x\) * Some covariate values
are only observed in control group $\to$ Cannot estimate
\(\mathbb{E}[Y(1) | X = x\)

Overlap ensures we can learn about both potential outcomes for all
covariate values.

\subsection{Identification Result}\label{_identification_result}

\begin{theorem}[Identification of ATE]
\label{thm:ate_identification}
Under Assumptions 1.1 (unconfoundedness) and 1.2 (overlap):
\[
\text{ATE} = \E_X \left[ \E[Y_i \mid T_i = 1, X_i] - \E[Y_i \mid T_i = 0, X_i] \right]
\]
\end{theorem}

\begin{proof}
By the law of iterated expectations:
\[
\E[Y_i(1)] = \E_X[\E[Y_i(1) \mid X_i]]
\]

By unconfoundedness:
\[
\E[Y_i(1) \mid X_i] = \E[Y_i(1) \mid T_i = 1, X_i]
\]

By definition of observed outcomes:
\[
\E[Y_i(1) \mid T_i = 1, X_i] = \E[Y_i \mid T_i = 1, X_i]
\]

Therefore:
\[
\E[Y_i(1)] = \E_X[\E[Y_i \mid T_i = 1, X_i]]
\]

Similarly for $\E[Y_i(0)]$. Taking the difference yields the result.
\end{proof}

\begin{remark}
The ATE can be expressed as a function of the \textbf{observed} data distribution $\Prob(Y, T, X)$, making it statistically estimable.
\end{remark}

\section{Deep Dive: Insurance Pricing
Example}\label{_deep_dive_insurance_pricing_example}

Let's work through the insurance pricing example in detail with concrete
numbers to build intuition for potential outcomes, confounding, and
identification.

\subsection{Setup}\label{_setup}

An annuity provider wants to estimate the causal effect of competitor
pricing on their weekly sales. They collect 52 weeks of data:

\begin{itemize}
\item
  \(Y_i\): Sales volume (number of annuities sold) in week \(i\)
\item
  \(T_i \in \{0, 1\}\): Competitor pricing indicator (1 = high prices, 0
  = low prices)
\item
  \(X_i = (X_{i1}, X_{i2}, X_{i3})\): Macroeconomic confounders

  \begin{itemize}
  \item
    \(X_{i1}\): VIX (market volatility index)
  \item
    \(X_{i2}\): Consumer sentiment index
  \item
    \(X_{i3}\): 10-year treasury rate
  \end{itemize}
\end{itemize}

\subsection{The Potential Outcomes}\label{_the_potential_outcomes}

For each week \(i\), two potential outcomes exist:

\begin{itemize}
\item
  \(Y_i(0)\): Sales if competitors had low prices that week
\item
  \(Y_i(1)\): Sales if competitors had high prices that week
\end{itemize}

\textbf{Week 1 example}: Suppose in reality \(T_1 = 1\) (competitors had
high prices) and we observed \(Y_1 = 245\) sales.

The potential outcomes are: * \(Y_1(1) = 245\): Observed (this actually
happened) * \(Y_1(0) = ?\): Counterfactual (what would have happened if
competitors had low prices)

We might guess \(Y_1(0) = 180\) (fewer sales with cheaper competitor
alternatives), implying an individual treatment effect
\(\tau_1 = 245 - 180 = 65\) additional sales from high competitor
prices.

\textbf{But this is fundamentally unknowable} for week 1 alone.

\subsection{Confounding in Action}\label{_confounding_in_action}

Why can't we just compare weeks with high vs. low competitor prices?

\textbf{Naive comparison} (biased):

\[\mathbb{E}[Y_i | T_i = 1] - \mathbb{E}[Y_i | T_i = 0]\]

This is biased because:

\begin{enumerate}
\def\labelenumi{\arabic{enumi}.}
\item
  \textbf{Economic cycles}: Competitors raise prices during high-demand
  periods (high VIX, low consumer sentiment)
\item
  \textbf{Selection}: Weeks with \(T_i = 1\) differ systematically from
  weeks with \(T_i = 0\)
\item
  \textbf{Confounding}: \(X_i\) affects both \(T_i\) and \(Y_i\)
\end{enumerate}

\textbf{Numerical example}:

Suppose the data shows: * Average sales when \(T_i = 1\):
\(\bar{Y}_{\text{high}} = 240\) * Average sales when \(T_i = 0\):
\(\bar{Y}_{\text{low}} = 190\) * Naive difference: \(240 - 190 = 50\)

But if we stratify by VIX:

\begin{longtable}[]{@{}
  >{\raggedright\arraybackslash}p{(\columnwidth - 6\tabcolsep) * \real{0.2500}}
  >{\raggedright\arraybackslash}p{(\columnwidth - 6\tabcolsep) * \real{0.2500}}
  >{\raggedright\arraybackslash}p{(\columnwidth - 6\tabcolsep) * \real{0.2500}}
  >{\raggedright\arraybackslash}p{(\columnwidth - 6\tabcolsep) * \real{0.2500}}@{}}
\caption{Sales by Competitor Price and VIX}\tabularnewline
\toprule\noalign{}
\begin{minipage}[b]{\linewidth}\raggedright
VIX Level
\end{minipage} & \begin{minipage}[b]{\linewidth}\raggedright
\(T=0\) avg sales
\end{minipage} & \begin{minipage}[b]{\linewidth}\raggedright
\(T=1\) avg sales
\end{minipage} & \begin{minipage}[b]{\linewidth}\raggedright
Difference
\end{minipage} \\
\midrule\noalign{}
\endfirsthead
\toprule\noalign{}
\begin{minipage}[b]{\linewidth}\raggedright
VIX Level
\end{minipage} & \begin{minipage}[b]{\linewidth}\raggedright
\(T=0\) avg sales
\end{minipage} & \begin{minipage}[b]{\linewidth}\raggedright
\(T=1\) avg sales
\end{minipage} & \begin{minipage}[b]{\linewidth}\raggedright
Difference
\end{minipage} \\
\midrule\noalign{}
\endhead
\bottomrule\noalign{}
\endlastfoot
Low (\textless{} 15) & 210 & 270 & +60 \\
High ($\geq$ 15) & 170 & 210 & +40 \\
\end{longtable}

\textbf{Key observation}: Within each VIX stratum, the treatment effect
is positive (high competitor prices $\to$ more sales). But:

\begin{itemize}
\item
  Competitors tend to raise prices when VIX is low (high demand periods)
\item
  Low VIX weeks have higher baseline sales regardless of competitor
  pricing
\item
  The naive comparison confounds the treatment effect with the VIX
  effect
\end{itemize}

\textbf{True ATE} (averaging within-stratum effects):
\(\frac{60 + 40}{2} = 50\)

In this simple example, the naive and conditional estimates coincide
\textbf{by accident}. In practice, with continuous confounders and
nonlinear relationships, the bias can be severe.

\subsection{Verifying Assumptions}\label{_verifying_assumptions}

\textbf{Unconfoundedness}: Is
\(\{Y_i(0), Y_i(1)\} \perp\!\!\!\perp T_i \mid X_i\) plausible?

\begin{itemize}
\item
  \textbf{Plausible}: If competitor pricing decisions are driven by
  observable macro conditions (VIX, sentiment, rates), then conditional
  on \(X_i\), the pricing is as-if random.
\item
  \textbf{Violation}: If competitors have private information about
  demand shocks (e.g., proprietary consumer surveys), then \(T_i\) is
  related to potential outcomes even after conditioning on \(X_i\).
\end{itemize}

\textbf{Overlap}: Is \(0 < \mathbb{P}(T_i = 1 | X_i) < 1\)?

\begin{itemize}
\item
  \textbf{Satisfied}: If there exist weeks with both high and low
  competitor prices across all VIX/sentiment/rate combinations
\item
  \textbf{Violated}: If competitors \textbf{never} raise prices when VIX
  \textgreater{} 25 $\to$ No treated units with high VIX $\to$ Cannot estimate
  \(\mathbb{E}[Y(1) | \text{VIX} > 25\)
\end{itemize}

We'll explore overlap violations in detail next.

\section{Example 2: Healthcare Treatment
Effects}\label{_example_2_healthcare_treatment_effects}

To reinforce the potential outcomes framework, let's examine a different
domain: estimating the effect of a medication on patient outcomes.

\subsection{Setup: Clinical Observational
Study}\label{_setup_clinical_observational_study}

A hospital wants to estimate the causal effect of a new diabetes
medication on HbA1c levels (blood sugar control). They have
observational data on 1,000 patients:

\begin{itemize}
\item
  \(Y_i\): Change in HbA1c after 6 months (negative = improvement)
\item
  \(T_i \in \{0, 1\}\): Treatment indicator (1 = new medication, 0 =
  standard medication)
\item
  \(X_i\): Patient characteristics

  \begin{itemize}
  \item
    Age, baseline HbA1c, BMI, comorbidities, insurance type, physician
    ID
  \end{itemize}
\end{itemize}

\textbf{Why observational?} The medication is already approved, so
randomization would be unethical. Doctors prescribe based on patient
characteristics.

\subsection{Potential Outcomes in
Healthcare}\label{_potential_outcomes_in_healthcare}

For each patient \(i\): * \(Y_i(1)\): Change in HbA1c if given new
medication * \(Y_i(0)\): Change in HbA1c if given standard medication

\textbf{Patient 42 example}: Suppose patient 42 receives the new
medication (\(T_{42} = 1\)) and experiences \(Y_{42} = -1.8\) (HbA1c
drops 1.8 points).

\begin{itemize}
\item
  \(Y_{42}(1) = -1.8\): Observed outcome (this actually happened)
\item
  \(Y_{42}(0) = ?\): Counterfactual (what would have happened with
  standard medication)
\end{itemize}

\textbf{Individual treatment effect}:
\(\tau_{42} = Y_{42}(1) - Y_{42}(0)\)

If we knew \(Y_{42}(0) = -1.0\), the individual effect would be
\(\tau_{42} = -1.8 - (-1.0) = -0.8\) (an additional 0.8 point reduction
from the new medication). But \(Y_{42}(0)\) is fundamentally
unobservable.

\subsection{Confounding by Indication}\label{_confounding_by_indication}

\textbf{Confounding by indication} is a classic problem in healthcare:
doctors prescribe treatments based on patient characteristics that also
affect outcomes.

\textbf{Scenario}: Suppose doctors prescribe the new medication
primarily to: * Younger patients (better adherence) * Patients with
higher baseline HbA1c (more room for improvement) * Patients with fewer
comorbidities (lower risk of side effects)

Now: * \(T_i\) is related to age, baseline HbA1c, comorbidities * These
variables also affect \(Y_i\) (younger patients may improve more
regardless of medication)

\textbf{Naive comparison} (biased):

\[\mathbb{E}[Y_i | T_i = 1] - \mathbb{E}[Y_i | T_i = 0]\]

might show the new medication is much better, but this confounds: *
\textbf{True medication effect}: \(\mathbb{E}[Y_i(1) - Y_i(0)\) *
\textbf{Patient selection}: Treated patients are younger, healthier,
more likely to improve anyway

\subsection{Unconfoundedness in
Healthcare}\label{_unconfoundedness_in_healthcare}

\textbf{Assumption}: Conditional on patient characteristics \(X_i\),
treatment assignment is "as-if randomized":

\[\{Y_i(0), Y_i(1)\} \perp\!\!\!\perp T_i \mid X_i\]

\textbf{When plausible}: * If doctors prescribe based solely on
observable characteristics (age, baseline HbA1c, comorbidities, etc.) *
All relevant patient features are recorded in electronic health records

\textbf{When violated}: * Doctors have private information not in the
EHR (patient motivation, family support, subtle clinical signs) *
Patients self-select into treatment based on unobservable factors (fear
of side effects, personal preferences)

\subsection{Overlap in Healthcare}\label{_overlap_in_healthcare}

\textbf{Overlap}: For all values of \(X\), some patients receive
treatment and some receive control:

\[0 < \mathbb{P}(T_i = 1 | X_i = x) < 1\]

\textbf{Example violation}: Suppose doctors \textbf{never} prescribe the
new medication to patients over age 75 (concern about kidney function).

\begin{itemize}
\item
  For \(\text{age} > 75\): \(e(x) = 0\) $\to$ No treated patients $\to$ Cannot
  estimate \(\mathbb{E}[Y(1) | \text{age} > 75\)
\item
  Cannot estimate treatment effect for elderly patients without
  extrapolation
\end{itemize}

\textbf{Practical solution}: Either: 1. \textbf{Restrict population}:
Estimate ATE only for ages 18-75 (where overlap holds) 2.
\textbf{Collect more data}: Find hospitals that do prescribe to elderly
patients 3. \textbf{Randomized trial}: If effect on elderly is crucial,
conduct RCT

\subsection{Why This Example Matters}\label{_why_this_example_matters}

Healthcare provides clear intuition for key concepts:

\begin{enumerate}
\def\labelenumi{\arabic{enumi}.}
\item
  \textbf{Individual effects unknowable}: Can't give patient 42 both
  medications simultaneously
\item
  \textbf{Confounding by indication}: Treatment decisions based on
  prognosis
\item
  \textbf{Unconfoundedness plausibility}: Depends on EHR completeness
  and physician decision-making
\item
  \textbf{Overlap violations}: Age limits, contraindications create
  deterministic rules
\item
  \textbf{Ethical constraints}: Observational data often necessary when
  randomization is unethical
\end{enumerate}

\textbf{Connection to insurance pricing}: Same framework, different
domain: * Healthcare: Doctors select patients for treatment based on
characteristics * Insurance: Competitors set prices based on market
conditions * Both: Need to condition on confounders to identify causal
effects

\section{Understanding Overlap: A Critical
Requirement}\label{_understanding_overlap_a_critical_requirement}

The overlap assumption is often overlooked but is critical for causal
inference. Let's see what happens when it fails.

\subsection{Overlap Defined}\label{_overlap_defined}

\textbf{Overlap (Positivity)}: For all \(x\) in the support of \(X\):

\[0 < e(x) < 1 \quad \text{where} \quad e(x) = \mathbb{P}(T_i = 1 | X_i = x)\]

The function \(e(x)\) is the \textbf{propensity score}---the probability
of treatment given confounders.

\textbf{Interpretation}: * \(e(x) > 0\): Some units with \(X = x\)
receive treatment $\to$ Can estimate \(\mathbb{E}[Y(1) | X = x\) *
\(e(x) < 1\): Some units with \(X = x\) receive control $\to$ Can estimate
\(\mathbb{E}[Y(0) | X = x\) * Both required to estimate
\(\mathbb{E}[Y(1) | X = x\) - \textbackslash mathbb\{E\}{[}Y(0)
\textbar{} X = x{]}{]}

\subsection{What Happens When Overlap
Fails?}\label{_what_happens_when_overlap_fails}

\textbf{Example: Extreme confounding}

Suppose competitor pricing depends deterministically on VIX:

\[T_i = \begin{cases}
1 & \text{if VIX}_i < 20 \\
0 & \text{if VIX}_i \geq 20
\end{cases}\]

Now: * For \(\text{VIX} < 20\): \(e(x) = 1\) $\to$ All weeks have high
competitor prices $\to$ Never observe \(Y_i(0)\) * For
\(\text{VIX} \geq 20\): \(e(x) = 0\) $\to$ All weeks have low competitor
prices $\to$ Never observe \(Y_i(1)\)

\textbf{Implication}: We cannot estimate
\(\mathbb{E}[Y(0) | \text{VIX} < 20\) or
\(\mathbb{E}[Y(1) | \text{VIX} \geq 20\).

The ATE is \textbf{not identified} without additional assumptions (e.g.,
parametric extrapolation).

\subsection{Near-Violations}\label{_near_violations}

In practice, overlap often holds technically but is weak:

\begin{itemize}
\item
  \(e(x) = 0.02\): Only 2\% of units with \(X = x\) are treated $\to$ Huge
  variance in \(\hat{\mathbb{E}}[Y(1) | X = x\)
\item
  \(e(x) = 0.98\): Only 2\% of units with \(X = x\) are controls $\to$ Huge
  variance in \(\hat{\mathbb{E}}[Y(0) | X = x\)
\end{itemize}

\textbf{Solution approaches}: 1. \textbf{Trimming}: Drop units with
\(e(x)\) close to 0 or 1 (changes estimand to trimmed population) 2.
\textbf{Weighting}: Propensity score weighting (upweight rare treatment
assignments) 3. \textbf{Regularization}: Shrink extreme propensity score
weights

Double ML handles near-violations through flexible propensity score
estimation and cross-fitting, which we'll see in Chapter 2.

\section{Python Implementation: Propensity Score
Methods}\label{_python_implementation_propensity_score_methods}

Let's demonstrate propensity score estimation and check overlap violations.

\begin{minted}{python}
import numpy as np
import pandas as pd
from sklearn.linear_model import LogisticRegression
from sklearn.ensemble import RandomForestClassifier
import matplotlib.pyplot as plt

np.random.seed(42)
n = 500

# Generate confounders
X = np.random.randn(n, 3)

# Generate propensity score with overlap
# e(X) = P(T=1|X) = logit^{-1}(X1 + 0.5*X2)
logit_e = X[:, 0] + 0.5 * X[:, 1]
e_true = 1 / (1 + np.exp(-logit_e))

# Generate treatment
T = np.random.binomial(1, e_true)

# Generate outcomes with treatment effect = 3
Y0 = X[:, 0] + X[:, 1]**2 + np.random.randn(n)
Y1 = Y0 + 3  # True ATE = 3
Y = T * Y1 + (1 - T) * Y0

print("=== Overlap Check ===")
print(f"Min propensity score: {e_true.min():.3f}")
print(f"Max propensity score: {e_true.max():.3f}")
print(f"Treatment prevalence: {T.mean():.3f}")
print(f"Overlap satisfied: {(e_true > 0.01).all() and (e_true < 0.99).all()}")

# Estimate propensity score with logistic regression
ps_model = LogisticRegression(penalty=None, max_iter=1000)
ps_model.fit(X, T)
e_hat = ps_model.predict_proba(X)[:, 1]

# Estimate propensity score with random forest (more flexible)
rf_model = RandomForestClassifier(n_estimators=100, max_depth=5, random_state=42)
rf_model.fit(X, T)
e_hat_rf = rf_model.predict_proba(X)[:, 1]

# Compare true vs. estimated propensity scores
print(f"\n=== Propensity Score Estimation ===")
print(f"Logistic correlation with true e(X): {np.corrcoef(e_true, e_hat)[0,1]:.3f}")
print(f"RF correlation with true e(X): {np.corrcoef(e_true, e_hat_rf)[0,1]:.3f}")

# Inverse Propensity Weighting (IPW) estimator
# ATE_IPW = E[T*Y/e(X)] - E[(1-T)*Y/(1-e(X))]
ate_ipw = (T * Y / e_hat).mean() - ((1 - T) * Y / (1 - e_hat)).mean()
print(f"\n=== ATE Estimation ===")
print(f"True ATE: 3.000")
print(f"Naive (ignore confounding): {Y[T==1].mean() - Y[T==0].mean():.3f}")
print(f"IPW (logistic e(X)): {ate_ipw:.3f}")

# Check for extreme weights (diagnostics)
weights_treated = 1 / e_hat[T == 1]
weights_control = 1 / (1 - e_hat[T == 0])
print(f"\n=== Weight Diagnostics ===")
print(f"Max treated weight: {weights_treated.max():.2f}")
print(f"Max control weight: {weights_control.max():.2f}")
print(f"Effective sample size (treated): {(weights_treated.sum())**2 / (weights_treated**2).sum():.0f}/{T.sum()}")
print(f"Effective sample size (control): {(weights_control.sum())**2 / (weights_control**2).sum():.0f}/{(1-T).sum()}")
\end{minted}

\textbf{Expected output}:

\begin{verbatim}
=== Overlap Check ===
Min propensity score: 0.076
Max propensity score: 0.924
Treatment prevalence: 0.500
Overlap satisfied: True

=== Propensity Score Estimation ===
Logistic correlation with true e(X): 0.974
RF correlation with true e(X): 0.956

=== ATE Estimation ===
True ATE: 3.000
Naive (ignore confounding): 0.234
IPW (logistic e(X)): 2.987

=== Weight Diagnostics ===
Max treated weight: 13.19
Max control weight: 13.16
Effective sample size (treated): 213/250
Effective sample size (control): 214/250
\end{verbatim}

\textbf{Key observations}:

\begin{enumerate}
\def\labelenumi{\arabic{enumi}.}
\item
  \textbf{Overlap satisfied}: All units have \(0.076 < e(X) < 0.924\)
\item
  \textbf{Naive estimator biased}: Ignoring confounders yields
  \(\hat{\text{ATE}} = 0.234\) vs. true 3.0
\item
  \textbf{IPW corrects bias}: Propensity score weighting recovers
  \(\hat{\text{ATE}} = 2.987 \approx 3.0\)
\item
  \textbf{Effective sample size}: Weighting reduces effective sample
  from 250 to \textasciitilde213 (15\% loss)
\end{enumerate}

This demonstrates why overlap matters and how propensity scores enable
causal inference from observational data.

\section{The Frisch-Waugh-Lovell
Theorem}\label{_the_frisch_waugh_lovell_theorem}

The Frisch-Waugh-Lovell (FWL) theorem
\cite{frisch1933partial,lovell1963seasonal} is a classical result in
regression analysis that provides geometric intuition for how Double ML
isolates treatment effects.

\subsection{Regression Setup}\label{_regression_setup}

Consider the linear regression:

\[Y_i = \beta_0 + \beta_1 T_i + \beta_2' X_i + \epsilon_i\]

where: * \(Y_i\): Outcome * \(T_i\): Treatment (scalar) * \(X_i\):
Confounders (vector) * \(\beta_1\): Treatment effect parameter of
interest

\textbf{Question}: Can we estimate \(\beta_1\) without explicitly
controlling for \(X_i\) in the same regression?

\subsection{The FWL Theorem}\label{_the_fwl_theorem}

\begin{theorem}[Frisch-Waugh-Lovell]
\label{thm:fwl}
The coefficient $\hat{\beta}_1$ from regressing $Y$ on $(T, X)$ is \textbf{identical} to the coefficient from the following two-step procedure:
\begin{enumerate}
\item \textbf{Residualize treatment}: Regress $T_i$ on $X_i$ to obtain residuals $\tilde{T}_i = T_i - \E[T_i \mid X_i]$
\item \textbf{Residualize outcome}: Regress $Y_i$ on $X_i$ to obtain residuals $\tilde{Y}_i = Y_i - \E[Y_i \mid X_i]$
\item \textbf{Regress residuals}: Regress $\tilde{Y}_i$ on $\tilde{T}_i$
\end{enumerate}

Formally:
\[
\hat{\beta}_1 = \frac{\Cov(\tilde{Y}_i, \tilde{T}_i)}{\Var(\tilde{T}_i)}
\]
\end{theorem}

\begin{proof}
The OLS estimator for $\beta_1$ in the full regression is:
\[
\hat{\beta}_1 = (T'M_X T)^{-1} T'M_X Y
\]
where $M_X = I - X(X'X)^{-1}X'$ is the projection matrix onto the space orthogonal to $X$.

Note that:
\begin{itemize}
\item $M_X T = T - X(X'X)^{-1}X'T = \tilde{T}$ (residuals from regressing $T$ on $X$)
\item $M_X Y = Y - X(X'X)^{-1}X'Y = \tilde{Y}$ (residuals from regressing $Y$ on $X$)
\end{itemize}

Therefore:
\[
\hat{\beta}_1 = (\tilde{T}'\tilde{T})^{-1} \tilde{T}'\tilde{Y} = \frac{\tilde{T}'\tilde{Y}}{\tilde{T}'\tilde{T}}
\]
which is the OLS coefficient from regressing $\tilde{Y}$ on $\tilde{T}$.
\end{proof}

\subsection{Geometric Interpretation}\label{_geometric_interpretation}

\textbf{Key insight}: The FWL theorem says that controlling for \(X\) is
equivalent to:

\begin{enumerate}
\def\labelenumi{\arabic{enumi}.}
\item
  \textbf{Removing the part of \(T\) explained by \(X\)}:
  \(\tilde{T}_i\) is the variation in treatment \textbf{orthogonal} to
  confounders
\item
  \textbf{Removing the part of \(Y\) explained by \(X\)}:
  \(\tilde{Y}_i\) is the variation in outcome \textbf{orthogonal} to
  confounders
\item
  \textbf{Regressing orthogonalized outcome on orthogonalized treatment}
\end{enumerate}

This is \textbf{partialling out}: we isolate the relationship between
\(T\) and \(Y\) that is not explained by \(X\).

\subsection{Mathematical Properties of
Residualization}\label{_mathematical_properties_of_residualization}

Residualization has important mathematical properties that explain why
it works for causal inference.

\begin{theorem}[Variance Reduction]
\label{thm:variance_reduction}
The residualized treatment has variance less than or equal to the original treatment:
\[
\Var(\tilde{T}_i) \leq \Var(T_i)
\]
with equality if and only if $T$ and $X$ are uncorrelated.
\end{theorem}

\begin{proof}
By definition of residuals:
\[
\tilde{T}_i = T_i - \E[T_i \mid X_i]
\]

The variance decomposition gives:
\[
\Var(T_i) = \Var(\E[T_i \mid X_i]) + \E[\Var(T_i \mid X_i)]
\]

But $\E[T_i \mid X_i]$ and $\tilde{T}_i$ are uncorrelated (by construction of residuals), so:
\[
\Var(T_i) = \Var(\E[T_i \mid X_i]) + \Var(\tilde{T}_i)
\]

Rearranging:
\[
\Var(\tilde{T}_i) = \Var(T_i) - \Var(\E[T_i \mid X_i])
\]

Since $\Var(\E[T_i \mid X_i]) \geq 0$, we have $\Var(\tilde{T}_i) \leq \Var(T_i)$.
\end{proof}

\begin{remark}
\textbf{Interpretation}:
\begin{itemize}
\item $\Var(\E[T_i \mid X_i])$: Variation in treatment \textbf{explained} by confounders
\item $\Var(\tilde{T}_i)$: Variation in treatment \textbf{unexplained} by confounders (the part we use for identification)
\end{itemize}

\textbf{Why this matters}: If confounders strongly predict treatment ($X$ explains most variation in $T$), then:
\begin{itemize}
\item $\Var(\tilde{T}_i)$ is small $\to$ Less ``identifying variation''
\item Standard errors on $\hat{\beta}_1$ are large $\to$ Imprecise estimates
\item This is the \textbf{strong confounding} problem---mitigated by larger sample sizes or stronger instruments
\end{itemize}
\end{remark}

\textbf{Connection to \(R^2\)}: Define the first-stage \(R^2\) as:

\[R^2_{T \sim X} = \frac{\text{Var}(\mathbb{E}[T_i | X_i])}{\text{Var}(T_i)} = 1 - \frac{\text{Var}(\tilde{T}_i)}{\text{Var}(T_i)}\]

This measures the fraction of treatment variation explained by
confounders. For causal inference: * \(R^2_{T \sim X} \approx 0\): Weak
confounding $\to$ Large identifying variation *
\(R^2_{T \sim X} \approx 1\): Strong confounding $\to$ Small identifying
variation $\to$ Imprecise estimates

\textbf{Example}: In our insurance pricing setup: * If VIX, sentiment,
and rates strongly predict competitor pricing (\(R^2 = 0.85\)), only
15\% of treatment variation is unexplained by confounders * This 15\%
residual variation is what identifies the treatment effect * Large
\(R^2\) is good for prediction, but makes causal inference harder

\subsection{Orthogonal Decomposition
(Advanced)}\label{_orthogonal_decomposition_advanced}

The FWL theorem is a special case of orthogonal decomposition in Hilbert
space. Let \(\mathcal{H}\) be the space of square-integrable random
variables with inner product
\(\langle Y, Z \rangle = \mathbb{E}[YZ\).

\textbf{Decomposition}: Any \(T \in \mathcal{H}\) can be uniquely
decomposed as:

\[T = P_X T + (I - P_X) T = \mathbb{E}[T | X] + \tilde{T}\]

where: * \(P_X\): Projection operator onto the subspace spanned by \(X\)
* \(\mathbb{E}[T | X\): Component of \(T\) in the \(X\)-subspace *
\(\tilde{T}\): Component of \(T\) orthogonal to \(X\)

\textbf{Orthogonality}: By definition of projection:

\[\langle \tilde{T}, h(X) \rangle = \mathbb{E}[\tilde{T} \cdot h(X)] = 0 \quad \text{for all} \quad h\]

This says \(\tilde{T}\) is uncorrelated with any function of \(X\).

\textbf{FWL in Hilbert space notation}:

\begin{align}
\beta_1 &= \frac{\langle Y, \tilde{T} \rangle}{\langle \tilde{T}, \tilde{T} \rangle} \\
&= \frac{\langle \tilde{Y}, \tilde{T} \rangle}{\langle \tilde{T}, \tilde{T} \rangle} \quad \text{(since } \langle Y - \tilde{Y}, \tilde{T} \rangle = 0\text{)}
\end{align}

This geometric view clarifies \textbf{why} FWL works: we're computing
the treatment effect using only the components of \(Y\) and \(T\) that
are orthogonal to confounders.

\subsection{Connection to Double Machine
Learning}\label{_connection_to_double_machine_learning}

FWL theorem motivates the Double ML approach
\cite{chernozhukov2018double}:

\begin{itemize}
\item
  \textbf{Classical FWL}: Uses \textbf{linear} regression to partial out
  \(X\)
\item
  \textbf{Double ML}: Uses \textbf{flexible machine learning} (random
  forests, boosting, neural networks) to partial out \(X\)
\end{itemize}

Why ML? When \(X\) is high-dimensional or the functional forms
\(\mathbb{E}[T|X]\) and \(\mathbb{E}[Y|X\) are nonlinear, linear
regression is misspecified. ML methods can approximate these conditional
expectations flexibly.

\textbf{Crucial addition}: Double ML adds sample splitting and
cross-fitting to avoid overfitting bias, which we'll develop in Chapter
2.

\subsection{Why Machine Learning? The Nonlinearity
Problem}\label{_why_machine_learning_the_nonlinearity_problem}

Linear regression for partialling out assumes \(\mathbb{E}[T|X]\) and
\(\mathbb{E}[Y|X\) are linear in \(X\). When this fails, FWL produces
biased estimates.

\textbf{Example: Nonlinear confounding}

Suppose the true data-generating process is:

\begin{align}
T_i &= X_{i1}^2 + X_{i2} + \epsilon_i^T \\
Y_i &= 2 T_i + X_{i1}^2 + X_{i2}^3 + \epsilon_i^Y
\end{align}

where: * \(\mathbb{E}[T|X]\) is quadratic in \(X_1\) *
\(\mathbb{E}[Y|X]\) is cubic in \(X_2\) * True treatment effect
\(\beta_1 = 2\)

\textbf{Linear FWL fails}: 1. Regress \(T\) on \(X\) (linear regression)
$\to$ Misspecified, underestimates \(\mathbb{E}[T|X]\) 2. Regress \(Y\) on
\(X\) (linear regression) $\to$ Misspecified, underestimates
\(\mathbb{E}[Y|X]\) 3. Residuals \(\tilde{T}, \tilde{Y}\) still
contain confounding from \(X\) 4. Final estimate
\(\hat{\beta}_1 \neq 2\) is biased

\textbf{Machine learning solution}: * Use \textbf{random forests} or
\textbf{boosting} to estimate \(\mathbb{E}[T|X]\) nonparametrically *
Use \textbf{neural networks} or \textbf{splines} to estimate
\(\mathbb{E}[Y|X]\) flexibly * Residuals properly remove nonlinear
confounding * Final estimate converges to \(\beta_1 = 2\)

\subsection{The High-Dimensional
Setting}\label{_the_high_dimensional_setting}

Modern applications often have: * \(p\) confounders where \(p\) is large
(hundreds or thousands) * \(p\) may even exceed \(n\) (more variables
than observations)

\textbf{Examples}: * \textbf{Healthcare}: Electronic health records with
thousands of diagnosis codes, lab values, medications *
\textbf{Marketing}: User demographics, browsing history, click patterns,
device info * \textbf{Insurance}: Competitor product features, market
conditions, economic indicators

\textbf{Linear regression breaks down}: * \(p > n\): Regression is
undefined (underdetermined system) * \(p \approx n\): Regression is
highly unstable (overfitting) * Large \(p\): Need regularization (Lasso,
Ridge) $\to$ Introduces bias

\textbf{Machine learning handles high dimensions}: * \textbf{Random
forests}: Split on important variables, ignore noise variables *
\textbf{Gradient boosting}: Sequentially add weak learners focusing on
residuals * \textbf{Neural networks}: Learn low-dimensional
representations * \textbf{Lasso regression}: Automatic variable
selection with \(\ell_1\) penalty

\subsection{The Regularization Bias
Problem}\label{subsec:ch01_regularization_bias}

When using regularized estimators (Lasso, Ridge, neural networks), a new
problem emerges: \textbf{regularization bias}.

\textbf{The issue}: Suppose we use Lasso to estimate
\(\hat{m}(X) = \mathbb{E}[Y|X\). The Lasso estimate is:

\[\hat{m}(X) = \arg\min_m \mathbb{E}[(Y - m(X))^2] + \lambda \|m\|_1\]

The penalty term \(\lambda \|m\|_1\) introduces bias:

\[\mathbb{E}[\hat{m}(X)] \neq \mathbb{E}[Y|X]\]

even with infinite data!

\textbf{Naive FWL with Lasso}: 1. Estimate \(\hat{m}_T(X)\) with Lasso $\to$
biased toward zero 2. Compute \(\tilde{T}_i = T_i - \hat{m}_T(X_i)\) $\to$
Residuals \textbf{not} mean-zero 3. Estimate \(\hat{\beta}_1\) $\to$ Biased
due to regularization bias in step 1

\textbf{Double ML solution (preview)}:

The key insight \cite{chernozhukov2018double} is that the treatment
effect \(\beta_1\) is \textbf{orthogonal} to the nuisance parameters
\(m_T(X), m_Y(X)\) in the sense that:

\[\frac{\partial}{\partial m} \mathbb{E}[(Y - \beta_1 T - m(X))^2] \bigg|_{m = m_0} = 0\]

This \textbf{Neyman orthogonality} means that small errors in
\(\hat{m}(X)\) don't affect \(\hat{\beta}_1\) to first order.

\textbf{Solution approach}: 1. \textbf{Sample splitting}: Use different
data for estimating \(\hat{m}(X)\) and \(\hat{\beta}_1\) 2.
\textbf{Cross-fitting}: Repeat with role reversal and average 3.
\textbf{Result}: Regularization bias cancels out,
\(\hat{\beta}_1 \to \beta_1\)

We'll develop this rigorously in Chapter 2 with the Neyman orthogonality
condition and the DML algorithm.

\subsection{When to Use Linear FWL vs. Double
ML}\label{_when_to_use_linear_fwl_vs_double_ml}

\textbf{Use linear FWL when}: * \(p\) is small (\textless{} 10
confounders) * Relationships are approximately linear * Sample size is
moderate (linear regression requires \(n > p\)) * Interpretability is
crucial (coefficients have direct meaning)

\textbf{Use Double ML when}: * \(p\) is large (hundreds or thousands of
confounders) * Relationships are nonlinear (interactions, polynomials,
thresholds) * Sample size is large (ML methods need data to learn
flexibly) * Prediction accuracy is more important than interpretability

\textbf{Example: Insurance pricing}: * \textbf{Linear FWL}: If only VIX,
sentiment, treasury rates (3 confounders) * \textbf{Double ML}: If using
hundreds of macro indicators, competitor product features, regional
demographics

\textbf{Example: Healthcare}: * \textbf{Linear FWL}: If only age,
baseline HbA1c, BMI (3 confounders) * \textbf{Double ML}: If using full
EHR (thousands of diagnosis codes, lab values, medications)

\textbf{Rule of thumb}: \(p > 20\) or suspected nonlinearity $\to$ Try
Double ML

\subsection{Computational
Considerations}\label{_computational_considerations}

\textbf{Linear FWL}: * Fast: \(O(np^2)\) for \(n\) observations, \(p\)
confounders * Scales to \(n = 10^6\), \(p = 10^3\) * No hyperparameter
tuning needed

\textbf{Double ML with Random Forests}: * Moderate:
\(O(n \log n \cdot p \cdot \text{trees})\) * Requires hyperparameter
tuning (max depth, min samples per leaf) * Parallelizes well (n\_jobs=48
on 64-core system) * Scales to \(n = 10^6\), \(p = 10^4\)

\textbf{Double ML with Neural Networks}: * Slow: Depends on architecture
and optimization * Requires careful hyperparameter tuning (layers,
width, learning rate, regularization) * Benefits from GPU acceleration *
Best for very large \(n\) (\(>10^6\)) and complex nonlinearity

\textbf{Practical workflow}: 1. Start with linear FWL (fast baseline) 2.
Try Random Forest DML (good default for nonlinearity) 3. Use neural
networks only if RF insufficient and computational budget allows

\section{Python Implementation: FWL
Theorem}\label{_python_implementation_fwl_theorem}

Let's demonstrate the FWL theorem with a simple simulation.

\begin{minted}{python}
import numpy as np
from sklearn.linear_model import LinearRegression

# Set seed for reproducibility
np.random.seed(42)
n = 1000

# Generate data
X = np.random.randn(n, 3)  # 3 confounders
T = X[:, 0] + 0.5 * X[:, 1] + np.random.randn(n)  # Treatment depends on X
Y = 2 * T + X[:, 1] - X[:, 2] + np.random.randn(n)  # True effect = 2

# Method 1: Full regression (Y ~ T + X)
X_T = np.column_stack([T, X])
reg_full = LinearRegression().fit(X_T, Y)
beta1_full = reg_full.coef_[0]

print(f"Method 1 (Full regression): $\beta$$_1$ = {beta1_full:.4f}")

# Method 2: FWL two-step procedure
# Step 1: Residualize T on X
reg_T = LinearRegression().fit(X, T)
T_resid = T - reg_T.predict(X)

# Step 2: Residualize Y on X
reg_Y = LinearRegression().fit(X, Y)
Y_resid = Y - reg_Y.predict(X)

# Step 3: Regress residuals
reg_resid = LinearRegression().fit(T_resid.reshape(-1, 1), Y_resid)
beta1_fwl = reg_resid.coef_[0]

print(f"Method 2 (FWL two-step): $\beta$$_1$ = {beta1_fwl:.4f}")
print(f"Difference: {abs(beta1_full - beta1_fwl):.2e}")
print(f"True effect: 2.0000")
\end{minted}

\textbf{Expected output}:

\begin{verbatim}
Method 1 (Full regression): $\beta$$_1$ = 1.9845
Method 2 (FWL two-step): $\beta$$_1$ = 1.9845
Difference: 0.00e+00
True effect: 2.0000
\end{verbatim}

\textbf{Key observation}: Both methods yield \textbf{identical}
estimates (up to numerical precision), confirming the FWL theorem.

\section{Summary}\label{sec:ch01_summary}

This chapter developed the foundations for causal inference and Double
Machine Learning:

\begin{enumerate}
\def\labelenumi{\arabic{enumi}.}
\item
  \textbf{Potential Outcomes Framework}: Defines causal effects through
  counterfactuals \(Y_i(1), Y_i(0)\)
\item
  \textbf{Fundamental Problem}: Individual effects
  \(\tau_i = Y_i(1) - Y_i(0)\) are never observed
\item
  \textbf{Average Treatment Effect}: Population-level effect
  \(\text{ATE} = \mathbb{E}[Y_i(1) - Y_i(0)\) is identifiable
\item
  \textbf{Identification}: Under unconfoundedness and overlap, ATE can
  be expressed using observed data
\item
  \textbf{FWL Theorem}: Partialling out confounders \(X\) from \(T\) and
  \(Y\) isolates the treatment effect
\end{enumerate}

\textbf{Next chapter}: We extend these ideas to \textbf{nonlinear}
partialling out using machine learning, introducing the Neyman
orthogonality condition and the DML algorithm.

\section{Concluding Remarks}\label{_concluding_remarks}

This chapter established the foundations for causal inference in
observational studies. Three key insights emerged:

\textbf{1. Causal inference is fundamentally a missing data problem}

We can never observe both \(Y_i(0)\) and \(Y_i(1)\) for the same unit.
Individual treatment effects \(\tau_i = Y_i(1) - Y_i(0)\) are logically
unobservable. This is not a statistical limitation---no amount of data
or sophisticated estimation can recover individual effects without
strong assumptions (e.g., time travel or parallel universes).

The resolution: Focus on \textbf{population-level} effects (ATE) that
can be identified under plausible assumptions (unconfoundedness +
overlap). This shift from individual to average effects is the core move
in modern causal inference.

\textbf{2. Identification requires assumptions, but they can be
empirically checked}

Unconfoundedness \(\{Y_i(0), Y_i(1)\} \perp\!\!\!\perp T_i \mid X_i\) is
never directly testable---it involves counterfactuals. However, its
plausibility can be assessed: * \textbf{Domain knowledge}: Do we believe
all confounders are observed? * \textbf{Sensitivity analysis}: How much
unobserved confounding would be needed to overturn conclusions? *
\textbf{Overlap diagnostics}: Are propensity scores well-behaved? (We
can test this!) * \textbf{Placebo tests}: Do we find effects where we
shouldn't? (Falsification)

Overlap \(0 < e(x) < 1\) \textbf{is} testable: We can directly examine
the empirical propensity score distribution and check for violations or
near-violations. Trimming or restricting the population to regions with
good overlap is often necessary.

\textbf{3. FWL is the bridge from linear regression to modern causal ML}

The Frisch-Waugh-Lovell theorem shows that "controlling for confounders"
is mathematically equivalent to: 1. Residualizing treatment on
confounders: \(\tilde{T} = T - \mathbb{E}[T|X\) 2. Residualizing
outcome on confounders: \(\tilde{Y} = Y - \mathbb{E}[Y|X\) 3.
Regressing orthogonalized outcome on orthogonalized treatment

This partialling-out interpretation generalizes beyond linear
regression: * \textbf{Classical econometrics}: Use linear regression for
\(\mathbb{E}[T|X]\) and \(\mathbb{E}[Y|X\) * \textbf{Modern ML}:
Use random forests, boosting, or neural networks for flexible
approximation

But naively replacing linear regression with ML introduces
\textbf{regularization bias}---penalized estimators (Lasso, Ridge,
neural nets) are biased even asymptotically. Double ML solves this
through Neyman orthogonality and cross-fitting.

\subsection{Roadmap to Chapter 2}\label{_roadmap_to_chapter_2}

Chapter 2 develops the Double Machine Learning framework rigorously:

\textbf{Neyman Orthogonality}: We'll show why the treatment effect
\(\beta_1\) is \textbf{orthogonal} to nuisance parameters
\(\mathbb{E}[T|X]\), \textbackslash mathbb\{E\}{[}Y\textbar X{]}{]} in a
precise sense. This orthogonality means regularization bias in nuisance
estimation doesn't contaminate \(\hat{\beta}_1\) to first order.

\textbf{The DML Algorithm}: The complete procedure with: 1.
\textbf{Sample splitting}: Divide data into \(K\) folds 2.
\textbf{Cross-fitting}: Estimate nuisance parameters on one fold,
treatment effect on another 3. \textbf{Aggregation}: Average across
folds

\textbf{Theoretical guarantees}: Under high-level conditions: *
\(\sqrt{n}\)-consistency and asymptotic normality * Valid confidence
intervals using cross-fit standard errors * Robustness to slow
convergence of ML estimators

\textbf{Python implementation}: Working code using EconML with: * Random
Forest nuisance estimation * Inference with confidence intervals *
Comparison to naive approaches

By the end of Chapter 2, you'll have a complete, validated DML
implementation ready for the insurance competitor pricing application in
Chapter 4.

\section{Exercises}\label{sec:ch01_exercises}

\subsection{Conceptual Problems}\label{_conceptual_problems}

\textbf{Exercise 1.1} (Potential Outcomes): In the insurance pricing
example, write out the potential outcomes \(Y_i(0), Y_i(1)\) explicitly
in words. What assumption would make them equal (no treatment effect)?

\textbf{Solution}: For week \(i\): * \(Y_i(1)\): Number of annuities
sold in week \(i\) if competitors had high prices * \(Y_i(0)\): Number
of annuities sold in week \(i\) if competitors had low prices

For no treatment effect: \(Y_i(1) = Y_i(0)\) for all \(i\). This
requires that competitor pricing has \textbf{zero causal effect} on our
sales. Plausible if: * Customers are completely insensitive to
competitor prices (no price shopping) * Our product is perfectly
differentiated (no substitution) * Competitor prices don't signal market
conditions

\textbf{Exercise 1.2} (Fundamental Problem): Explain why collecting more
data does not solve the fundamental problem of causal inference. What
would we need to observe to compute individual treatment effects?

\textbf{Solution}: The fundamental problem is that we can never observe
\(Y_i(1)\) and \(Y_i(0)\) for the same unit \(i\) simultaneously---this
is logically impossible. Collecting more data gives us: * More weeks
with \(T_i = 1\) (high competitor prices) * More weeks with \(T_i = 0\)
(low competitor prices)

But for any specific week \(i\), we still only observe one potential
outcome. To compute \(\tau_i = Y_i(1) - Y_i(0)\), we would need to: 1.
\textbf{Time travel}: Observe week \(i\) twice with different competitor
prices, or 2. \textbf{Parallel universes}: Observe both treatments
simultaneously in the same week

Neither is feasible. We can only estimate \textbf{average} effects
across units.

\textbf{Exercise 1.3} (Unconfoundedness Violation): Give an example
where unconfoundedness is violated in the insurance pricing setting.
What variable might be unobserved that affects both competitor pricing
and sales?

\textbf{Solution}: Suppose competitors have access to a proprietary
consumer confidence survey (not publicly available) that predicts
insurance demand. They use this to set prices: * High survey confidence
$\to$ Competitors raise prices (high demand expected) * Low survey
confidence $\to$ Competitors lower prices (low demand expected)

This survey also affects our sales: * High confidence $\to$ More consumers
buy (including from us) * Low confidence $\to$ Fewer consumers buy

Now unconfoundedness is violated: * \(T_i\) (competitor pricing) is
affected by unobserved confidence * \(Y_i(0), Y_i(1)\) (our sales under
both treatments) are also affected by unobserved confidence *
Conditioning on observed \(X_i\) (VIX, sentiment, treasury rates) is
insufficient

Result: \(\{Y_i(0), Y_i(1)\} \not\perp\!\!\!\perp T_i \mid X_i\)

\textbf{Exercise 1.4} (Overlap Implications): Suppose overlap is
violated: for VIX \textgreater{} 25, we never observe high competitor
prices (\(e(x) = 0\) for VIX \textgreater{} 25). Can we still estimate:
(a) The ATE for the full population? (b) The ATE conditional on VIX
\textless{} 25?

\textbf{Solution}:

(a) \textbf{No}. The full population ATE is:

\[\text{ATE} = \mathbb{E}_X[\mathbb{E}[Y(1) | X] - \mathbb{E}[Y(0) | X]]\]

For VIX \textgreater{} 25, we cannot estimate
\(\mathbb{E}[Y(1) | \text{VIX} > 25\) (no treated units with high
VIX). Without this, the full ATE is not identified.

(b) \textbf{Yes}. The conditional ATE is:

\[\text{ATE}_{\text{VIX} < 25} = \mathbb{E}[\tau_i | \text{VIX} < 25]\]

Within the VIX \textless{} 25 subpopulation, overlap holds (both high
and low competitor prices observed). We can estimate both
\(\mathbb{E}[Y(1) | X, \text{VIX} < 25\) and
\(\mathbb{E}[Y(0) | X, \text{VIX} < 25\).

\textbf{Caveat}: This changes the estimand from "ATE for all weeks" to
"ATE for low-VIX weeks".

\subsection{Mathematical Problems}\label{_mathematical_problems}

\textbf{Exercise 1.5} (Variance Decomposition): Prove that
\(\text{Var}(\tilde{T}_i) \leq \text{Var}(T_i)\) where
\(\tilde{T}_i = T_i - \mathbb{E}[T_i | X_i\).

\textbf{Solution}: (Already proven in Theorem 1.3, reproduced here for
completeness)

By the law of total variance:

\[\text{Var}(T_i) = \mathbb{E}[\text{Var}(T_i | X_i)] + \text{Var}(\mathbb{E}[T_i | X_i])\]

By construction, \(\tilde{T}_i\) and \(\mathbb{E}[T_i | X_i\) are
uncorrelated:

\[\mathbb{E}[\tilde{T}_i \cdot \mathbb{E}[T_i | X_i]] = \mathbb{E}[\mathbb{E}[\tilde{T}_i | X_i] \cdot \mathbb{E}[T_i | X_i]] = 0\]

since \(\mathbb{E}[\tilde{T}_i | X_i\) = 0{]} by definition of
residuals.

Therefore:

\[\text{Var}(T_i) = \text{Var}(\mathbb{E}[T_i | X_i]) + \text{Var}(\tilde{T}_i)\]

Rearranging:

\[\text{Var}(\tilde{T}_i) = \text{Var}(T_i) - \text{Var}(\mathbb{E}[T_i | X_i]) \leq \text{Var}(T_i)\]

since variances are non-negative. Equality holds iff
\(\text{Var}(\mathbb{E}[T_i | X_i\)) = 0 \textbackslash iff T\_i
\textbackslash perp\textbackslash!\textbackslash!\textbackslash!\textbackslash perp
X\_i{]}.

\textbf{Interpretation}: Residualization removes the variation in \(T\)
explained by \(X\), leaving only unexplained variation. This unexplained
variation is the "identifying variation" for the treatment effect.

\textbf{Exercise 1.6} (Propensity Score Bounds): Show that under
unconfoundedness, the ATE can be written as:

\[\text{ATE} = \mathbb{E}\left[\frac{T_i Y_i}{e(X_i)} - \frac{(1 - T_i) Y_i}{1 - e(X_i)}\right]\]

where \(e(X_i) = \mathbb{P}(T_i = 1 | X_i)\) is the propensity score.

\textbf{Solution}:

Start with the identification formula from Theorem 1.1:

\[\mathbb{E}[Y_i(1)] = \mathbb{E}_X[\mathbb{E}[Y_i | T_i = 1, X_i]]\]

By the law of iterated expectations:

\[\mathbb{E}[Y_i(1)] = \mathbb{E}\left[\mathbb{E}\left[\frac{T_i Y_i}{e(X_i)} \bigg| X_i\right]\right]\]

The inner expectation is:

\[\mathbb{E}\left[\frac{T_i Y_i}{e(X_i)} \bigg| X_i\right] = \frac{\mathbb{E}[T_i Y_i | X_i]}{e(X_i)} = \frac{e(X_i) \mathbb{E}[Y_i | T_i = 1, X_i]}{e(X_i)} = \mathbb{E}[Y_i | T_i = 1, X_i]\]

Therefore:

\[\mathbb{E}[Y_i(1)] = \mathbb{E}\left[\frac{T_i Y_i}{e(X_i)}\right]\]

Similarly:

\[\mathbb{E}[Y_i(0)] = \mathbb{E}\left[\frac{(1 - T_i) Y_i}{1 - e(X_i)}\right]\]

Taking the difference yields the Inverse Propensity Weighting (IPW)
estimator.

\subsection{Computational Problems}\label{_computational_problems}

\textbf{Exercise 1.7} (Nonlinear FWL Failure): Modify the FWL Python code to use nonlinear confounding:

\begin{minted}{python}
# Generate data with nonlinearity
X = np.random.randn(n, 2)
T = X[:, 0]**2 + X[:, 1] + np.random.randn(n)  # Quadratic in X1
Y = 2 * T + X[:, 0]**2 + X[:, 1]**3 + np.random.randn(n)  # True effect = 2
\end{minted}

Does linear FWL recover the true effect \(\beta_1 = 2\)? If not, what
estimator would work?

\textbf{Solution}: Linear FWL will be biased because: 1. Regressing
\(T\) on \(X\) (linear) cannot capture \(\mathbb{E}[T|X] = X_1^2 + X_2\)
2. Regressing \(Y\) on \(X\) (linear) cannot capture
\(\mathbb{E}[Y|X] = 2\mathbb{E}[T|X] + X_1^2 + X_2^3\)
3. Residuals \(\tilde{T}, \tilde{Y}\) still
contain nonlinear confounding

\textbf{What would work}: * \textbf{Polynomial regression}: Include
\(X_1^2, X_2^3\) in the regression * \textbf{Random Forest}: Use
RandomForestRegressor to estimate \(\mathbb{E}[T|X]\) and
\(\mathbb{E}[Y|X\) * \textbf{Double ML}: Use flexible ML in
cross-fitting framework (Chapter 2)

\textbf{Exercise 1.8} (Effective Sample Size): In the propensity score
example, we computed effective sample sizes of 213/250 for treated and
214/250 for controls. Explain why weighting reduces effective sample
size and when this is problematic.

\textbf{Solution}:

\textbf{Why weighting reduces effective sample size}:

Inverse propensity weights are \(w_i = 1/e(X_i)\) for treated and
\(w_i = 1/(1-e(X_i))\) for controls. When \(e(X_i)\) varies across
units: * Units with \(e(X_i) \approx 1\): Low weights (already likely to
be treated) * Units with \(e(X_i) \approx 0\): High weights (unlikely to
be treated but got treated anyway)

High variance in weights $\to$ Effective sample size \textless{} actual
sample size.

Formally:

\[n_{\text{eff}} = \frac{(\sum w_i)^2}{\sum w_i^2}\]

\textbf{When problematic}: * \textbf{Near-violations}: If some
\(e(X_i) \approx 0.01\), weights can be 100+, drastically reducing
\(n_{\text{eff}}\) * \textbf{Precision loss}: Effective 213/250 means we
lost 37 units worth of information (15\%) * \textbf{Variance inflation}:
Standard errors increase by \(\sqrt{250/213} \approx 1.08\)

\textbf{Remedies}: 1. \textbf{Trimming}: Drop units with
\(e(X_i) < 0.05\) or \(e(X_i) > 0.95\) 2. \textbf{Weight truncation}:
Cap weights at percentiles (e.g., 99th) 3. \textbf{Doubly robust
estimation}: Combine weighting with outcome modeling

\section{References}\label{_references}

bibliography::{[}{]}

